% ACL 2026 style / REALM@EMNLP 2026 (archival long, max 8p content)
% Style files: https://github.com/acl-org/acl-style-files (master)
% Compile with: pdflatex -> bibtex -> pdflatex x2
% Overleaf: set compiler to pdfLaTeX
\documentclass[]{article}
\usepackage[review]{acl}   % 'review' removes author info for blind review
\usepackage{times}
\usepackage{latexsym}
\usepackage{amsmath}
\usepackage{amssymb}
\usepackage{booktabs}
\usepackage{multirow}
\usepackage{graphicx}
\usepackage{microtype}
\usepackage{xcolor}
\usepackage[utf8]{inputenc}
\usepackage[T1]{fontenc}

% ---------------------------------------------------------------------------
\title{Stage-Wise Persona Attribution in Deep Research Agent Pipelines}

% TODO: add authors after blind review
% \author{...}

\begin{document}
\maketitle

% ---------------------------------------------------------------------------
\begin{abstract}
% TODO: 결과 나오면 완성
Deep Research Agents (DRAs) generate long-form reports through a multi-stage pipeline---planning, search, compression, and writing---conditioned on user personas.
A key open question is \emph{which} stage of this pipeline actually reflects persona conditioning, and how reliably.
We introduce \textbf{N-way stage-wise persona attribution}: given a report artifact produced by one of $N$ candidate personas, identify the ground-truth persona from intermediate pipeline outputs alone.
Applying this framework to 120 reports generated by a Gemini-2.5-Flash DRA backbone over PDR-Bench (25 personas, 5 domains), we find that [RESULT PLACEHOLDER].
We further show that shortcut controls (identifier masking, actionable-only conditioning) [RESULT PLACEHOLDER].
Our artifacts, tracing adapter, and evaluation code are released to facilitate reproducibility.
\end{abstract}

% ---------------------------------------------------------------------------
\section{Introduction}
\label{sec:intro}
% TODO: hook + gap + contributions
\textcolor{red}{[TODO: write after \S2 gap is confirmed]}

% Placeholder contributions list:
\begin{itemize}
  \item We define the \emph{N-way stage-wise attribution} task for DRA pipelines and provide a formal treatment of stage-wise accuracy trajectories.
  \item We build DRATracer, a deterministic LangGraph v2 event listener that captures four stage artifacts from a single pipeline run without post-hoc reconstruction.
  \item We conduct the first controlled attribution study on PDR-Bench, with hard-negative candidate sets and lexical-leakage shortcut controls.
  \item [RESULT: key finding about which stage]
\end{itemize}

% ---------------------------------------------------------------------------
\section{Related Work}
\label{sec:related}

\subsection{Personalized Deep Research}

Personalization in information retrieval has a long history grounded in user modeling and collaborative filtering~\cite{placeholder}.
Recent work has extended personalization to large language model (LLM) generation, with benchmarks such as LaMP~\cite{placeholder} testing whether models can condition outputs on user-specific history.
\textbf{PDR-Bench}~\cite{pdr-bench} is the first benchmark specifically targeting personalized \emph{deep research}---multi-step, long-form report generation conditioned on rich user personas.
PDR-Bench provides 25 personas and 250 English query pairs spanning five domains, with personas encoding both identity attributes (demographics, occupation) and actionable preferences (information-seeking style, decision-making patterns).

A key gap in PDR-Bench and prior personalization benchmarks is the absence of a \emph{causal} diagnostic: they measure whether a report \emph{looks} personalized, but not \emph{which pipeline stage} drove the personalization.
We address this gap by framing personalization as an N-way attribution problem applied at each pipeline stage independently.

\subsection{Authorship Attribution and Persona Identification}

Classical authorship attribution identifies the author of a text from a closed candidate set, relying on stylometric features such as function-word distributions and syntactic patterns~\cite{placeholder}.
Neural approaches have shifted toward dense embeddings and fine-tuned classifiers~\cite{placeholder}.

\textbf{AuthorMap}~\cite{authormap} is the closest prior work to ours: it frames persona attribution as an N-way classification problem, using an LLM to identify which persona a summary was written for.
However, AuthorMap operates on \textbf{pairwise} comparisons of one-shot summaries---a single-stage, single-output setting.
Our work departs from AuthorMap in two fundamental ways:
(1)~\textbf{N-way simultaneous} attribution ($N=3$ by default, with $N\in\{2,5\}$ as sensitivity analyses), without decomposing into pairwise decisions; and
(2)~\textbf{stage-wise} attribution across a four-stage DRA pipeline, enabling \emph{localization} of where in the pipeline the conditioning signal is expressed or lost.

\subsection{Multi-Step Agent Pipelines and Preference Fidelity}

Deep research agents such as \texttt{open\_deep\_research}~\cite{open-deep-research} implement multi-step LangGraph pipelines: a query is first elaborated into a research brief, then dispatched to parallel researcher subgraphs that search and compress evidence, before a final synthesis pass generates the report.
This pipeline architecture means that a persona conditioning signal applied once at input must propagate through four qualitatively distinct transformation stages.

\textbf{PrefDisco}~\cite{prefdisco} finds that preference accuracy degrades specifically in the reasoning and content-selection layers of multi-step pipelines, motivating a stage-wise evaluation rather than end-to-end accuracy alone.
We operationalize this motivation with a tracing adapter that captures deterministic artifacts at each stage, enabling per-stage attribution without post-hoc reconstruction.

% ---------------------------------------------------------------------------
\section{Method}
\label{sec:method}

\subsection{Task Formulation}

Let $P$ be a DRA pipeline and $\mathcal{C} = \{p_1, \dots, p_N\}$ be a candidate persona set ($N=3$ in our primary experiments).
$P$ is conditioned on one \textbf{ground-truth persona} $p^* \in \mathcal{C}$ and a research query $q$, and produces a sequence of \textbf{stage artifacts}:
\begin{equation}
  \mathcal{A}(q, p^*) = \bigl(a_{\text{plan}},\; a_{\text{search}},\; a_{\text{compress}},\; a_{\text{write}}\bigr).
\end{equation}
The \textbf{N-way stage-wise attribution} task is defined as: for each stage artifact $a \in \mathcal{A}$, identify $p^*$ from $\mathcal{C}$:
\begin{equation}
  \hat{p} = f\bigl(a,\; \mathcal{C}\bigr), \qquad \hat{p} \in \mathcal{C}.
\end{equation}
Attribution accuracy at stage $s$ is $\mathrm{Acc}_s = \mathbb{E}[\mathbf{1}[\hat{p}_s = p^*]]$, with chance level $1/N$.
Stage-wise accuracy profiles localize \emph{where} in the pipeline the conditioning signal is expressed: a flat profile suggests the signal never propagates; a peak at $a_{\text{plan}}$ that decays by $a_{\text{write}}$ suggests early capture but downstream dilution.

\subsection{DRA Pipeline and Artifact Schema}

We use \textbf{open\_deep\_research}~\cite{open-deep-research} as the DRA backbone, instantiated with Gemini~2.5~Flash.
The pipeline implements a four-stage LangGraph graph; Table~\ref{tab:artifacts} summarizes each stage's artifact.

\begin{table}[t]
\centering
\small
\begin{tabular}{llp{4.2cm}}
\toprule
\textbf{Stage} & \textbf{Artifact} & \textbf{Content} \\
\midrule
Planning    & $a_{\text{plan}}$     & Research brief: scoped topics, intended angle \\
Search      & $a_{\text{search}}$   & Queries + top-3 Serper organic results each \\
Compression & $a_{\text{compress}}$ & Per-topic synthesis of retrieved evidence \\
Writing     & $a_{\text{write}}$    & Full long-form research report \\
\bottomrule
\end{tabular}
\caption{Four stage artifacts captured per DRA run.}
\label{tab:artifacts}
\end{table}

\textbf{DRATracer} is our custom LangGraph v2 event listener that captures all four artifacts deterministically from a single pipeline run.
The key implementation challenge is that three researcher subgraphs run concurrently under the supervisor; DRATracer resolves per-artifact topic identity using parent run-ID correlation across \texttt{on\_chain\_start} and \texttt{on\_chain\_end} events.
All artifacts are serialized to a versioned JSON schema (\texttt{dra\_trace\_v2}) with SHA-256 prompt hashes and a query/source ledger for reconciliation.

\subsection{Attribution Protocol}

\paragraph{LLM Matcher.}
We use Upstage Solar-pro3 as the primary attribution model (temperature~$=0$, structured output).
Given artifact text $a$ and persona set $\mathcal{C}$ (order-shuffled to prevent position bias), the matcher returns a single predicted persona $\hat{p}$.
A secondary matcher (GPT-5.4-nano) is applied to the reference configuration (seed~0 only) as a replication check.

\paragraph{Non-LLM Baselines.}
We include three baselines: (1)~\emph{Random}, uniform over $\mathcal{C}$ (expected accuracy~$= 1/N$); (2)~\emph{BM25}, Okapi BM25 similarity between artifact and persona text; (3)~\emph{Embedding}, cosine similarity via a sentence encoder.

\paragraph{Hard Negatives.}
For each ground-truth persona, the two distractor personas are the top-2 same-domain personas by actionable-token Jaccard similarity ($\geq 0.2$ = hard negative; otherwise nearest fallback, seeded per task).
This ensures chance-level performance is a meaningful baseline.

\subsection{Shortcut Controls}
\label{sec:shortcut}

High attribution accuracy could arise from \textbf{lexical leakage}---name tokens, demographic phrases, or verbatim persona text appearing in the report---rather than genuine content-intent alignment.
We control for this with three generation-time conditions and one post-processing condition (Table~\ref{tab:shortcut}).

\begin{table}[t]
\centering
\small
\begin{tabular}{lp{2.8cm}l}
\toprule
\textbf{Condition} & \textbf{Manipulation} & \textbf{Tests} \\
\midrule
Full (baseline)       & Unmodified persona                          & Upper bound \\
Actionable-only       & Remove identity tokens                       & Identity leak? \\
Identity-only         & Keep only demographic tokens                 & Beyond demographics? \\
Identifier-masked     & NER removal from report text                 & Surface mentions? \\
\bottomrule
\end{tabular}
\caption{Shortcut control conditions. Generation-time conditions modify the conditioning persona; identifier masking is a post-processing pass on existing artifacts.}
\label{tab:shortcut}
\end{table}

% ---------------------------------------------------------------------------
\section{Experimental Setup}
\label{sec:setup}

\subsection{Datasets}

\paragraph{PDR-Bench (Primary).}
We use the English split of PDR-Bench~\cite{pdr-bench}, which provides 25 user personas and 250 query-persona pairs across five domains (Education, Career, Health, Finance, Creative).
We apply domain-stratified largest-remainder sampling to select 20 confirmatory query groups; eligibility filtering removes tasks with fewer than three personas that have at least one actionable token after identity masking.
Hard-negative candidate sets are constructed as described in \S\ref{sec:method}.
The confirmatory run uses two seeds ($s \in \{0, 1\}$), yielding \textbf{120 reports} (20~groups~$\times$~3~personas~$\times$~2~seeds).
A 5-group development split (15~reports) serves as an engineering gate and is excluded from all research-question analyses.

\paragraph{LaMP-QA (Replication).}
We sample 15 queries from LaMP-QA~\cite{lamp-qa} across three categories: Art \& Entertainment, Lifestyle, and Society \& Culture (5 per category, stratified).
Distractors are the top-2 Jaccard matches from a 500-item same-category pool.
With two seeds: \textbf{90 reports}.
We treat LaMP-QA as a replication dataset only.

\subsection{Generation Stack}

All reports are generated with open\_deep\_research~\cite{open-deep-research} using Gemini~2.5~Flash for all pipeline roles.
Search is provided by a Serper adapter (max 7 queries per report, top-3 organic results each, max 15 unique URLs, global hard stop at 2,400 queries).
Generation seed is fixed via \texttt{generation\_config.seed}; other hyperparameters use model defaults.
Every artifact records a full \texttt{execution\_config} block including model IDs, seeds, and SHA-256 hashes for reproducibility.

\subsection{Report Budget}

Table~\ref{tab:budget} summarizes all generation runs.
Identifier-masked variants are produced by post-processing existing artifacts at zero additional cost.

\begin{table}[t]
\centering
\small
\begin{tabular}{lrrrr}
\toprule
\textbf{Run} & \textbf{Groups} & \textbf{$N$} & \textbf{Seeds} & \textbf{Reports} \\
\midrule
PDR-Bench confirmatory (full)      & 20 & 3 & 2 & 120 \\
PDR ablation---actionable-only     &  5 & 3 & 2 &  30 \\
PDR ablation---identity-only       &  5 & 3 & 2 &  30 \\
PDR ablation---shuffled-actionable &  5 & 3 & 1 &  15 \\
LaMP-QA replication                & 15 & 3 & 2 &  90 \\
\midrule
\textbf{Total}                     &    &   &   & \textbf{285} \\
\bottomrule
\end{tabular}
\caption{Generation budget (excludes 15-report dev pilot).}
\label{tab:budget}
\end{table}

\subsection{Evaluation Metrics}

\paragraph{Attribution Accuracy (Acc@1).}
Fraction of reports for which the matcher's top-1 prediction equals $p^*$.
Reported per stage and macro-averaged across seeds. Chance level: $1/N = 0.33$.

\paragraph{Stage Trajectory.}
The four-tuple $(\mathrm{Acc}_{\text{plan}}, \mathrm{Acc}_{\text{search}}, \mathrm{Acc}_{\text{compress}}, \mathrm{Acc}_{\text{write}})$ plotted over pipeline depth.
A rising trajectory implies signal \emph{accumulation}; a falling trajectory implies \emph{early capture but downstream dilution}.

\paragraph{Seed Variance.}
Std.\ dev.\ of Acc@1 across seeds 0 and 1 as a fragility indicator.

\paragraph{Inter-Matcher Agreement.}
Cohen's $\kappa$ between Solar-pro3 and GPT-5.4-nano on seed-0 PDR-Bench confirmatory split.

% ---------------------------------------------------------------------------
\section{Results}
\label{sec:results}
\textcolor{red}{[TODO: fill after batch runs complete]}

% ---------------------------------------------------------------------------
\section{Analysis}
\label{sec:analysis}
\textcolor{red}{[TODO: shortcut ablation, stage trajectory discussion]}

% ---------------------------------------------------------------------------
\section{Conclusion}
\label{sec:conclusion}
\textcolor{red}{[TODO]}

% ---------------------------------------------------------------------------
\bibliography{refs}

\end{document}
